% !TEX program = xelatex
\documentclass[11pt,a4paper]{ctexart}
\usepackage{geometry}
\geometry{left=2.5cm, right=2.5cm, top=3cm, bottom=3cm}
\usepackage{listings}
\usepackage{xcolor}
\usepackage{hyperref}
\usepackage{fancyhdr}
\usepackage{graphicx}
\usepackage{tikz} 
\usepackage{linegoal}
\usepackage{calligra} 
\usepackage{amsmath}
\usepackage{amssymb}

\hypersetup{colorlinks=true, linkcolor=black, urlcolor=blue}
\setcounter{tocdepth}{2}
\pagestyle{fancy}
\fancyhf{}
\fancyhead[L]{算法竞赛核心模板库}
\fancyhead[R]{\thepage}
\setlength{\headheight}{15pt}

\lstset{
    language=C++,
    basicstyle=\ttfamily\small,
    keywordstyle=\color{blue}\bfseries,
    commentstyle=\color[rgb]{0,0.6,0},
    stringstyle=\color[rgb]{0.58,0,0.82},
    numbers=left,
    numberstyle=\tiny\color{gray},
    breaklines=true,
    frame=single,
    tabsize=4,
    showstringspaces=false,
    extendedchars=false,
}

\begin{document}

\begin{titlepage}
    \centering
    \vspace*{2cm} 

    \begin{tikzpicture}
        \node [circle, minimum size=6.5cm, path picture={
            \node at (path picture bounding box.center){
                \includegraphics[width=7cm]{cover.png}
            };
        }] {};
    \end{tikzpicture}

    \vspace{2cm} 

    {\huge \textbf{\kaishu 算法竞赛核心模板库} \par}
    \vspace{0.5cm}
    {\color{gray!60} \rule{0.6\textwidth}{0.5pt} \vspace{2pt} \par}
    {\color{gray!60} \rule{0.6\textwidth}{0.5pt} \par}

    \vspace{2cm} 
    {\Large \hspace{-1cm} {\calligra \huge Jiangcheneee} \hspace{1cm} \Large \fangsong 华中农业大学 \par}
    \vspace{0.8cm} 
    {\color{gray!60} \rule{2cm}{0.5pt} \par}
    \vspace{0.8cm} 
    {\large \the\year 年 \the\month 月 \the\day 日 \par}
    \vfill 
\end{titlepage}

\thispagestyle{empty}
\newpage
\tableofcontents
\newpage
\section{一键启动模板}
\subsection{启动文件}
\lstinputlisting{"启动文件.cpp"}

\section{stl库}
\subsection{algorithm}
\lstinputlisting{"stl库/algorithm.cpp"}

\subsection{atan2}
\lstinputlisting{"stl库/atan2.cpp"}

\subsection{bitset}
\lstinputlisting{"stl库/bitset.cpp"}

\subsection{map}
\lstinputlisting{"stl库/map.cpp"}

\subsection{pbds\_tree}
\lstinputlisting{"stl库/pbds_tree.cpp"}

\subsection{priority\_queue}
\lstinputlisting{"stl库/priority_queue.cpp"}

\subsection{queue}
\lstinputlisting{"stl库/queue.cpp"}

\subsection{set\&\&multiset}
\lstinputlisting{"stl库/set&&multiset.cpp"}

\subsection{stack}
\lstinputlisting{"stl库/stack.cpp"}

\subsection{string}
\lstinputlisting{"stl库/string.cpp"}

\subsection{tuple}
\lstinputlisting{"stl库/tuple.cpp"}

\subsection{vector}
\lstinputlisting{"stl库/vector.cpp"}

\section{博弈论}
\subsection{sg函数}
\lstinputlisting{"博弈论/sg函数.cpp"}

\section{图论}
\subsection{dij}
\lstinputlisting{"图论/dij.cpp"}

\subsection{欧拉路径}
\lstinputlisting{"图论/欧拉路径.cpp"}

\section{数学}
\subsection{FFT}
\lstinputlisting{"数学/FFT.cpp"}

\subsection{NTT}
\lstinputlisting{"数学/NTT.cpp"}

\subsection{exgcd}
\lstinputlisting{"数学/exgcd.cpp"}

\subsection{中国剩余定理}
\lstinputlisting{"数学/中国剩余定理.cpp"}

\subsection{中国剩余定理（推导笔记）}
\input{"数学/中国剩余定理（推导笔记）.tex"}

\subsection{任意模数NTT}
\lstinputlisting{"数学/任意模数NTT.cpp"}

\subsection{前缀基}
\lstinputlisting{"数学/前缀基.cpp"}

\subsection{前缀线性基}
\lstinputlisting{"数学/前缀线性基.cpp"}

\subsection{卢卡斯定理}
\lstinputlisting{"数学/卢卡斯定理.cpp"}

\subsection{卢卡斯定理（推导笔记）}
\input{"数学/卢卡斯定理（推导笔记）.tex"}

\subsection{多项式卷积（推导笔记）}
\input{"数学/多项式卷积（推导笔记）.tex"}

\subsection{欧拉筛}
\lstinputlisting{"数学/欧拉筛.cpp"}

\subsection{矩阵快速幂}
\lstinputlisting{"数学/矩阵快速幂.cpp"}

\subsection{矩阵快速幂（加速DP取max）}
\lstinputlisting{"数学/矩阵快速幂（加速DP取max）.cpp"}

\subsection{线性基（方案还原与删除）}
\lstinputlisting{"数学/线性基（方案还原与删除）.cpp"}

\subsection{线性基（贪心法）}
\lstinputlisting{"数学/线性基（贪心法）.cpp"}

\subsection{范德蒙德卷积（推导笔记）}
\input{"数学/范德蒙德卷积（推导笔记）.tex"}

\subsection{莫比乌斯反演（二维GCD求和）}
\lstinputlisting{"数学/莫比乌斯反演（二维GCD求和）.cpp"}

\subsection{莫比乌斯反演（推导笔记）}
\input{"数学/莫比乌斯反演（推导笔记）.tex"}

\subsection{莫比乌斯反演（特定GCD计数）}
\lstinputlisting{"数学/莫比乌斯反演（特定GCD计数）.cpp"}

\subsection{行列式计算}
\lstinputlisting{"数学/行列式计算.cpp"}

\section{数据结构}
\subsection{01trie}
\lstinputlisting{"数据结构/01trie.cpp"}

\subsection{dsu}
\lstinputlisting{"数据结构/dsu.cpp"}

\subsection{主席树（区间MEX）}
\lstinputlisting{"数据结构/主席树（区间MEX）.cpp"}

\subsection{主席树（区间元素数）}
\lstinputlisting{"数据结构/主席树（区间元素数）.cpp"}

\subsection{主席树（区间第K小）}
\lstinputlisting{"数据结构/主席树（区间第K小）.cpp"}

\subsection{区间mex}
\lstinputlisting{"数据结构/区间mex.cpp"}

\subsection{树状数组}
\lstinputlisting{"数据结构/树状数组.cpp"}

\subsection{线段树}
\lstinputlisting{"数据结构/线段树.cpp"}

\subsection{线段树（区间max）}
\lstinputlisting{"数据结构/线段树（区间max）.cpp"}

\subsection{莫队}
\lstinputlisting{"数据结构/莫队.cpp"}

\section{树论}
\subsection{倍增}
\lstinputlisting{"树论/倍增.cpp"}

\subsection{启发式合并}
\lstinputlisting{"树论/启发式合并.cpp"}

\section{计算几何}
\subsection{凸包（单调链法）}
\lstinputlisting{"计算几何/凸包（单调链法）.cpp"}

\subsection{极角排序}
\lstinputlisting{"计算几何/极角排序.cpp"}

\subsection{求凸包直径（旋转卡壳）}
\lstinputlisting{"计算几何/求凸包直径（旋转卡壳）.cpp"}

\subsection{求凸包面积}
\lstinputlisting{"计算几何/求凸包面积.cpp"}

\subsection{闵可夫斯基和与皮克定理}
\input{"计算几何/闵可夫斯基和与皮克定理.tex"}

\section{杂项}
\subsection{CDQ分治}
\lstinputlisting{"杂项/CDQ分治.cpp"}

\subsection{CDQ分治优化DP}
\lstinputlisting{"杂项/CDQ分治优化DP.cpp"}

\subsection{CDQ分治（推导笔记）}
\input{"杂项/CDQ分治（推导笔记）.tex"}

\subsection{int128}
\lstinputlisting{"杂项/int128.cpp"}

\subsection{modint}
\lstinputlisting{"杂项/modint.cpp"}


\end{document}
